\textbf{Phần 1. Trắc nghiệm lựa chọn}
\setcounterpageref{ex}{0}

\begin{ex} %[4D10V55]
Khi nhân cả tử số và mẫu số của một phân số với cùng một số tự nhiên khác 0 thì ta được
\choice
{\True Phân số bằng phân số ban đầu}
{Phân số lớn hơn phân số ban đầu}
{Phân số nhỏ hơn phân số ban đầu}
{Số tự nhiên}
\loigiai{
Được một phân số bằng phân số đã cho.\\
\textbf{Đáp án của câu hỏi là A}
}
\end{ex}

\begin{ex} %[4D10V55]
Khi chia cả tử số và mẫu số của một phân số cho cùng một số tự nhiên khác 0 (và cùng chia hết) ta được
\choice
{Phân số gấp đôi}
{Phân số nhỏ hơn}
{\True Phân số bằng phân số ban đầu}
{Phân số nghịch đảo}
\loigiai{
Được một phân số bằng phân số đã cho.\\
\textbf{Đáp án của câu hỏi là C}
}
\end{ex}

\begin{ex} %[4D10V55]
Điền số thích hợp vào chỗ chấm: $\dfrac{2}{3} = \dfrac{\dots}{9}$.
\choice
{$4$}
{\True $6$}
{$8$}
{$5$}
\loigiai{
Nhân cả tử và mẫu với 3: $2 \times 3 = 6$.\\
\textbf{Đáp án của câu hỏi là B}
}
\end{ex}

\begin{ex} %[4D10V55]
Điền số thích hợp vào chỗ chấm: $\dfrac{15}{20} = \dfrac{3}{\dots}$.
\choice
{$3$}
{$5$}
{$6$}
{\True $4$}
\loigiai{
Chia cả tử và mẫu cho 5: $20 : 5 = 4$.\\
\textbf{Đáp án của câu hỏi là D}
}
\end{ex}

\begin{ex} %[4D10V55]
Phân số nào sau đây bằng phân số $\dfrac{3}{5}$?
\choice
{\True $\dfrac{9}{15}$}
{$\dfrac{6}{15}$}
{$\dfrac{5}{3}$}
{$\dfrac{6}{8}$}
\loigiai{
$\dfrac{3 \times 3}{5 \times 3} = \dfrac{9}{15}$.\\
\textbf{Đáp án của câu hỏi là A}
}
\end{ex}

\begin{ex} %[4D10V55]
Phân số nào sau đây bằng phân số $\dfrac{12}{18}$?
\choice
{$\dfrac{3}{4}$}
{$\dfrac{4}{5}$}
{\True $\dfrac{2}{3}$}
{$\dfrac{1}{2}$}
\loigiai{
$\dfrac{12 : 6}{18 : 6} = \dfrac{2}{3}$.\\
\textbf{Đáp án của câu hỏi là C}
}
\end{ex}

\begin{ex} %[4D10V55]
Tìm $x$ biết: $\dfrac{x}{8} = \dfrac{15}{24}$.
\choice
{$3$}
{\True $5$}
{$6$}
{$4$}
\loigiai{
$24 : 3 = 8 \Rightarrow x = 15 : 3 = 5$.\\
\textbf{Đáp án của câu hỏi là B}
}
\end{ex}

\begin{ex} %[4D10V55]
Tìm $y$ biết: $\dfrac{4}{y} = \dfrac{28}{35}$.
\choice
{$4$}
{$6$}
{$7$}
{\True $5$}
\loigiai{
$28 : 7 = 4 \Rightarrow y = 35 : 7 = 5$.\\
\textbf{Đáp án của câu hỏi là D}
}
\end{ex}

\begin{ex} %[4D10V55]
Trong các phân số sau, phân số nào KHÔNG bằng $\dfrac{4}{6}$?
\choice
{\True $\dfrac{8}{10}$}
{$\dfrac{2}{3}$}
{$\dfrac{8}{12}$}
{$\dfrac{12}{18}$}
\loigiai{
$\dfrac{8}{10} = \dfrac{4}{5} \neq \dfrac{4}{6}$.\\
\textbf{Đáp án của câu hỏi là A}
}
\end{ex}

\begin{ex} %[4D10V55]
Viết phân số $\dfrac{3}{4}$ thành phân số có mẫu số bằng 100:
\choice
{$\dfrac{30}{100}$}
{$\dfrac{80}{100}$}
{\True $\dfrac{75}{100}$}
{$\dfrac{70}{100}$}
\loigiai{
$\dfrac{3 \times 25}{4 \times 25} = \dfrac{75}{100}$.\\
\textbf{Đáp án của câu hỏi là C}
}
\end{ex}

\begin{ex} %[4D10V55]
Viết phân số $\dfrac{1}{5}$ thành phân số có mẫu số là 20:
\choice
{$\dfrac{2}{20}$}
{\True $\dfrac{4}{20}$}
{$\dfrac{5}{20}$}
{$\dfrac{3}{20}$}
\loigiai{
$\dfrac{1 \times 4}{5 \times 4} = \dfrac{4}{20}$.\\
\textbf{Đáp án của câu hỏi là B}
}
\end{ex}

\begin{ex} %[4D10V55]
Cho phân số $\dfrac{7}{13}$. Hỏi phải cộng thêm vào cả tử số và mẫu số cùng một số tự nhiên nào để được phân số bằng $\dfrac{2}{3}$?
\choice
{$3$}
{$4$}
{$6$}
{\True $5$}
\loigiai{
Hiệu mẫu và tử: $13 - 7 = 6$.\\
Phân số mới có hiệu mẫu và tử là $3 - 2 = 1$ phần $= 6$.\\
Tử số mới: $2 \times 6 = 12 \Rightarrow$ Số cần cộng: $12 - 7 = 5$.\\
\textbf{Đáp án của câu hỏi là D}
}
\end{ex}

\textbf{Phần 2. Trắc nghiệm đúng sai}
\setcounterpageref{ex}{0}

\begin{ex} %[4D10V55]
Xét tính đúng sai của các khẳng định về tính chất cơ bản phân số:
\choiceTF
{\True $\dfrac{a}{b} = \dfrac{a \times m}{b \times m}$ ($m$ là số tự nhiên khác 0).}
{\True $\dfrac{a}{b} = \dfrac{a : n}{b : n}$ ($n$ là ước chung khác 0 của $a$ và $b$).}
{$\dfrac{a}{b} = \dfrac{a + m}{b + m}$.}
{\True $\dfrac{3}{4} = \dfrac{6}{8} = \dfrac{9}{12}$.}
\loigiai{
\begin{itemize}
    \item Nhân cùng 1 số: Đúng.
    \item Chia cùng 1 số: Đúng.
    \item Cộng cùng 1 số thì phân số thay đổi giá trị (không bằng nhau). Sai.
    \item Dãy phân số bằng nhau đúng.
\end{itemize}
\begin{itemize}
    \item a) \textbf{Đúng}.
    \item b) \textbf{Đúng}.
    \item c) \textbf{Sai}.
    \item d) \textbf{Đúng}.
\end{itemize}
}
\end{ex}

\begin{ex} %[4D10V55]
Cho phân số $\dfrac{18}{24}$. Xét tính đúng sai:
\choiceTF
{\True Rút gọn chia cho 6 được phân số $\dfrac{3}{4}$.}
{\True Phân số $\dfrac{18}{24}$ bằng phân số $\dfrac{9}{12}$.}
{Phân số $\dfrac{18}{24}$ bằng phân số $\dfrac{6}{7}$.}
{\True Nhân cả tử và mẫu với 2 được phân số $\dfrac{36}{48}$.}
\loigiai{
\begin{itemize}
    \item $18:6 / 24:6 = 3/4$. Đúng.
    \item $18:2 / 24:2 = 9/12$. Đúng.
    \item $3/4 \neq 6/7$. Sai.
    \item $18\times 2 / 24\times 2 = 36/48$. Đúng.
\end{itemize}
\begin{itemize}
    \item a) \textbf{Đúng}.
    \item b) \textbf{Đúng}.
    \item c) \textbf{Sai}.
    \item d) \textbf{Đúng}.
\end{itemize}
}
\end{ex}

\textbf{Phần 3. Trả lời ngắn}
\setcounterpageref{ex}{0}

\begin{ex} %[4D10V55]
Điền số thích hợp vào chỗ chấm: $\dfrac{3}{5} = \dfrac{15}{\dots}$.
\shortans{25}
\loigiai{
\textbf{Đáp số:} $25$\\
$5 \times 5 = 25$.
}
\end{ex}

\begin{ex} %[4D10V55]
Điền số thích hợp vào chỗ chấm: $\dfrac{24}{36} = \dfrac{\dots}{3}$.
\shortans{2}
\loigiai{
\textbf{Đáp số:} $2$\\
$36 : 12 = 3 \Rightarrow 24 : 12 = 2$.
}
\end{ex}

\begin{ex} %[4D10V55]
Tìm $x$ biết: $\dfrac{x}{12} = \dfrac{3}{4}$.
\shortans{9}
\loigiai{
\textbf{Đáp số:} $9$\\
$x = 3 \times 3 = 9$.
}
\end{ex}

\begin{ex} %[4D10V55]
Tìm $y$ biết: $\dfrac{35}{y} = \dfrac{5}{7}$.
\shortans{49}
\loigiai{
\textbf{Đáp số:} $49$\\
$y = 7 \times 7 = 49$.
}
\end{ex}

\begin{ex} %[4D10V55]
Viết phân số $\dfrac{2}{5}$ thành phân số có mẫu số bằng 50. Tìm tử số của phân số mới.
\shortans{20}
\loigiai{
\textbf{Đáp số:} $20$\\
$2 \times 10 = 20$.
}
\end{ex}

\begin{ex} %[4D10V55]
Tìm một phân số bằng $\dfrac{3}{4}$ biết tổng của tử số và mẫu số bằng 35. Tìm tử số của phân số đó.
\shortans{15}
\loigiai{
\textbf{Đáp số:} $15$\\
Tổng số phần: $3 + 4 = 7$ phần. Giá trị 1 phần: $35 : 7 = 5$.\\
Tử số là: $5 \times 3 = 15$.
}
\end{ex}
